\section{总结与延伸}

\subsection{核心贡献}

Anthropic 这篇工程博客的核心贡献在于提供了一套经过实践验证的、解决 Agent 跨 context window 工作问题的完整方案。其关键要素可以概括为：

\begin{enumerate}
    \item \textbf{双阶段架构}：Initializer Agent 搭建环境，Coding Agent 增量推进
    \item \textbf{Feature List}：结构化的 JSON 格式功能清单，防止 Agent 过早宣布完工
    \item \textbf{Progress File + Git History}：跨 session 的状态传递机制
    \item \textbf{标准化启动流程}：减少每个 session 的"热身"成本
    \item \textbf{端到端测试}：通过浏览器自动化确保功能真正可用
\end{enumerate}

\subsection{对 Agent 开发者的启示}

对于正在构建长时间运行 Agent 系统的开发者而言，本文最重要的启示是：

\begin{importantbox}{不要只优化模型，更要优化 Harness}
Agent 的表现不仅取决于底层模型的能力，更取决于围绕模型构建的 harness（运行框架）的设计。即使是最强的前沿模型，如果 harness 设计不当，也无法在长时间运行场景下保持有效输出。反之，一个精心设计的 harness 可以让模型的能力得到充分发挥。
\end{importantbox}

\subsection{拓展阅读}

\begin{itemize}
    \item Anthropic Claude Agent SDK 文档
    \item Anthropic Claude 4 Prompting Guide——关于多 context window 工作流的最佳实践
    \item 原文附带的 quickstart 代码示例
    \item Anthropic 关于 Agent 安全性的研究（长时间自主运行的 Agent 需要额外的安全考量）
\end{itemize}

%% --- Fin del contenido --- %%

\end{document}
